
\documentclass[11pt]{article}

\usepackage[a4paper,margin=1in]{geometry}
\usepackage{amsmath,amssymb}
\usepackage{siunitx}
\usepackage{graphicx}
\usepackage{hyperref}
\usepackage{bm}
\usepackage{physics}

\title{Momentum-Driven Self-Starting Siphons:\\
A Simple One-Dimensional Model and Its Limits}
\author{}
\date{\today}

\begin{document}

\maketitle

\begin{abstract}
We study self-starting siphons whose priming is driven by the momentum of the inflowing liquid rather than by external actions such as manual filling or pumping.
Starting from a minimal geometric description of the pipe and container, we develop a hierarchy of one-dimensional models.
We first derive an ideal lossless model with constant container level, which provides an upper bound on the achievable flow rates and crest heights.
We then include viscous friction and minor losses, and couple the siphon dynamics to the evolving container free-surface height.
The central focus is the start-up (priming) phase: we formulate a momentum-based condition under which the flow traverses the highest crest and a continuous liquid column is established.
We next examine the draining phase and characterize when the siphon naturally stops.
Finally, we delineate the main limits of the model (including surface-tension and cavitation constraints) and propose table-top experiments to validate and refine the theory.
\end{abstract}

\section{Introduction}

A siphon is a hydraulic device that transports liquid from a higher reservoir to a lower outlet through a tube whose highest point lies above the initial free surface.
Classical textbook treatments focus on siphons that have already been \emph{primed}---that is, the tube is completely filled with liquid and a continuous column connects inlet and outlet.
In practice, however, many practical and natural configurations involve \emph{self-starting} siphons, in which the tube is initially unprimed and the system must spontaneously develop a continuous liquid column without external pumping.

The aim of this paper is twofold:
\begin{enumerate}
  \item to construct a simple but explicit one-dimensional (1D) model that captures both the start-up and draining phases of a self-starting siphon, and
  \item to identify the main physical parameters that govern the maximum attainable crest height and the robustness of priming, together with the limits beyond which the model breaks down.
\end{enumerate}

To do so, we introduce a clear geometric parametrization of the siphon, derive governing equations from control-volume balances, and then progressively relax idealizations: we move from a constant-level, lossless model to models including Darcy--Weisbach friction, minor losses, and a time-varying container free surface.
Throughout, we emphasize the connection between the equations and physical intuition.

\section{Physical Picture and Flow Regimes}
\label{sec:physical-picture}

\subsection{Conventional (Non Self-Starting) Siphons}

In a conventional siphon, the tube is initially filled with liquid by external means.
This may be done by:
\begin{itemize}
  \item manually filling the tube with liquid,
  \item using a pump or syringe to evacuate air, or
  \item submerging the tube entirely and then positioning it so that its ends lie in the desired locations.
\end{itemize}
In these cases, the initial condition is a fully primed siphon: the tube contains a continuous liquid column, and the subsequent flow is driven entirely by gravity and pressure differences.
The dynamics are then well described by steady Bernoulli arguments augmented with head losses.

\subsection{Momentum-Driven Self-Starting Siphons}

Here we are interested in a different scenario: the siphon starts from an unprimed state.
The liquid level in the container rises, or the tube is inserted progressively, until the liquid first reaches a lower crest of the tube.
At that instant, the geometry of the tube is fixed and the container level is no longer externally controlled.
From that moment onward, the evolution is governed solely by gravity and fluid dynamics.

A key feature of \emph{momentum-driven} self-starting siphons is that, when the liquid first reaches the lower crest, it can already possess a nonzero downward velocity.
If the available head is sufficiently large, and the initial momentum is not dissipated too quickly by viscous and minor losses, the liquid can traverse the higher crest (or apex) and establish a continuous column.
The siphon then transitions to a fully primed state without external intervention.

We call the period before the geometry and water level are ``frozen'' the \emph{externally driven phase}, and the subsequent evolution the \emph{natural dynamics phase}.

\subsection{Regime Classification}

It is useful to organize siphon behavior into qualitative regimes:
\begin{itemize}
  \item \textbf{Externally primed siphons:} conventional textbook case; no start-up dynamics.
  \item \textbf{Momentum-driven self-starting siphons:} priming occurs spontaneously due to inertial effects as the liquid crosses one or more crests.
  \item \textbf{Failed priming:} the liquid reaches an intermediate crest but comes to rest before surmounting the highest point of the tube; the column breaks and no sustained flow is established.
  \item \textbf{Breakdown regimes:} even if primed, the siphon may fail due to cavitation, surface-tension effects, or strong air entrainment.
\end{itemize}
The following sections provide a quantitative framework to distinguish these regimes.

\section{Geometry and Coordinate System}
\label{sec:geometry}

We now define the geometry and the notation used throughout.
Consider a container with a free-surface height $H(t)$ above a fixed reference level (e.g.\ the container bottom).
The siphon is represented by a tube of constant inner diameter $D$ and total arc length $L_{\text{tot}}$, whose centerline is described by an arc-length coordinate $L \in [0,L_{\text{tot}}]$.

We denote by $h(L)$ the vertical height of the tube centerline at position $L$ above the same reference level.
Important points along the tube are:
\begin{itemize}
  \item Inlet point $A$ at $L = 0$ with height $h_A = h(0)$;
  \item Lower crest $B$ at $L = L_B$ with height $h_B$;
  \item Apex (highest point) $D$ at $L = L_D$ with height $h_D$;
  \item Outlet point $E$ at $L = L_{\text{tot}}$ with height $h_E$.
\end{itemize}

For simplicity, we assume the tube can be decomposed into straight segments and short curved segments (bends).
On each straight segment $i$, $h(L)$ is linear,
\begin{equation}
  h(L) = h_{i,0} + s_i (L - L_{i,0}), \qquad L \in [L_{i,0}, L_{i,1}],
\end{equation}
where $s_i$ is the constant slope of segment $i$.
Curved segments are treated either geometrically, via circular arcs, or---more simply---via minor-loss coefficients in the head-loss model (Section~\ref{sec:friction-losses}).

Throughout the main text we adopt a piecewise-linear approximation for $h(L)$, and account for bends through minor-loss coefficients.
A more detailed geometric treatment can be relegated to an appendix if needed.

\section{Governing Equations}
\label{sec:governing-equations}

We model the flow as incompressible and one-dimensional along the tube, with cross-sectional area
\begin{equation}
  A = \frac{\pi D^2}{4}.
\end{equation}
The volumetric flow rate is $Q(t)$ (assumed uniform along $L$), so the average velocity is $v(t) = Q(t)/A$.

\subsection{Ideal Constant-Level, Lossless Model}
\label{sec:ideal-model}

As a useful reference, consider an idealized configuration in which:
\begin{itemize}
  \item the container free surface remains at a fixed height $H_0$;
  \item the fluid is inviscid and there are no head losses; and
  \item the siphon is already primed and the flow is steady.
\end{itemize}

Applying Bernoulli's equation between the free surface and the outlet (neglecting velocity at the free surface) gives
\begin{equation}
  \frac{p_{\text{atm}}}{\rho g} + H_0 = \frac{p_{\text{atm}}}{\rho g} + h_E + \frac{v^2}{2g}.
\end{equation}
Thus the ideal outlet velocity is
\begin{equation}
  v_{\text{ideal}} = \sqrt{2g (H_0 - h_E)}, \qquad
  Q_{\text{ideal}} = A v_{\text{ideal}}.
\end{equation}
This provides an \emph{upper bound} on the discharge and on the dynamic pressure that can be used to overcome crests and losses.
Any realistic model must produce $|Q(t)| \le Q_{\text{ideal}}$ for the same head difference.

\subsection{Model with Friction and Minor Losses}
\label{sec:friction-losses}

We now include viscous and geometric losses.
For a fully developed, single-phase flow with average velocity $v(t)$, the Darcy--Weisbach formula gives the distributed head loss
\begin{equation}
  h_f(t) = f \,\frac{L_{\text{tot}}}{D} \,\frac{v(t)^2}{2g},
\end{equation}
where $f$ is the Darcy friction factor (which may depend on Reynolds number and roughness).
Additionally, we introduce a total minor-loss coefficient $K_m$ to account for bends, entrance, and exit effects, leading to an additional head loss
\begin{equation}
  h_m(t) = K_m \,\frac{v(t)^2}{2g}.
\end{equation}
The combined head loss is
\begin{equation}
  h_{\ell}(t) = h_f(t) + h_m(t)
  = \left( f\,\frac{L_{\text{tot}}}{D} + K_m \right)\frac{v(t)^2}{2g}.
\end{equation}

\subsubsection*{Tank Coupling and Unsteady Momentum Balance}

Let $A_{\text{tank}}$ be the cross-sectional area of the container.
Mass conservation gives
\begin{equation}
  \frac{dH}{dt} = -\,\frac{Q(t)}{A_{\text{tank}}}.
\end{equation}

A simple unsteady 1D momentum balance on the liquid column of length $L_{\text{tot}}$ can be written as
\begin{equation}
  \rho L_{\text{tot}} \frac{d}{dt}\!\left(\frac{Q}{A}\right)
  = \rho g A \Delta h_{\text{drive}}(t) - \rho g A\, h_{\ell}(t),
  \label{eq:unsteady-momentum}
\end{equation}
where $\Delta h_{\text{drive}}(t)$ is the effective driving head between the free surface at height $H(t)$ and the outlet at $h_E$,
\begin{equation}
  \Delta h_{\text{drive}}(t) \approx H(t) - h_E.
\end{equation}
Neglecting convective acceleration along the tube (rigid-column approximation) and dividing by $\rho A$, we obtain an ordinary differential equation (ODE) for $Q(t)$:
\begin{equation}
  \frac{L_{\text{tot}}}{A} \frac{dQ}{dt}
  = g A \bigl(H(t) - h_E\bigr)
  - g A \left( f\,\frac{L_{\text{tot}}}{D} + K_m \right)\frac{Q(t)^2}{2gA^2}.
\end{equation}
Equivalently,
\begin{equation}
  \frac{dQ}{dt}
  = \frac{gA^2}{L_{\text{tot}}}\bigl(H(t) - h_E\bigr)
  - \frac{A}{L_{\text{tot}}}\left( f\,\frac{L_{\text{tot}}}{D} + K_m \right)\frac{Q(t)^2}{2A^2}.
  \label{eq:Q-ODE}
\end{equation}
Together with the tank-level equation, this gives a closed 2D dynamical system for $(H(t), Q(t))$ during the draining phase, provided the siphon remains primed.

\section{Start-Up and Priming Dynamics}
\label{sec:priming}

We now focus on the moment when the liquid first reaches the lower crest $B$.
Before that instant, the container level and/or tube position may be controlled externally.
At the instant of interest:
\begin{itemize}
  \item the tube geometry is fixed;
  \item the container free surface height is $H_0$; and
  \item the liquid at $B$ has some velocity $v_0$ and may still contain small air pockets.
\end{itemize}

We idealize the liquid column between $B$ and the apex $D$ as a control volume.
To establish a continuous column across $D$, the liquid must be able to ascend from $h_B$ to $h_D$ despite losses along the path.
A simple energetic priming criterion is that the available head at $B$ exceeds the sum of elevation gain and head losses up to $D$:
\begin{equation}
  \frac{v_0^2}{2g}
  + \bigl(H_0 - h_B\bigr)
  \;\ge\;
  \bigl(h_D - h_B\bigr)
  + h_{f,\text{up}}(Q_0)
  + h_{m,\text{up}}(Q_0),
  \label{eq:priming-criterion}
\end{equation}
where $Q_0 = A v_0$ and $h_{f,\text{up}}, h_{m,\text{up}}$ denote the frictional and minor-loss contributions along the \emph{upstream} segment from $B$ to $D$.
If \eqref{eq:priming-criterion} is satisfied with some margin, the liquid can traverse the apex with nonzero velocity and begin to descend on the downstream side.
Once the downstream leg fills, the siphon becomes fully primed and the dynamics follow the ODEs of Section~\ref{sec:governing-equations}.

The criterion \eqref{eq:priming-criterion} makes clear how several parameters influence priming:
\begin{itemize}
  \item deeper submergence of the inlet (or of an intermediate point $C$) increases $H_0 - h_B$ and thus the available driving head;
  \item higher apex height $h_D$ makes priming more difficult;
  \item larger tube diameter $D$ reduces frictional losses but also changes $v_0$ for a given $Q_0$;
  \item increased roughness or sharper bends (larger $f$ and $K_m$) raise the required head for priming.
\end{itemize}

\section{Draining Phase and Shutdown}
\label{sec:draining}

Once the siphon is primed, the evolution of $H(t)$ and $Q(t)$ is governed by the coupled ODEs for tank level and flow rate.
As the container drains, the driving head $\Delta h_{\text{drive}}(t) = H(t) - h_E$ decreases.
The siphon stops when one of the following occurs:
\begin{enumerate}
  \item \textbf{Geometric shutdown:} $H(t)$ approaches $h_E$, so the hydrostatic head vanishes and the flow rate tends to zero.
  \item \textbf{Cavitation-induced shutdown:} the pressure at the apex $D$ falls below the vapor pressure of the liquid, causing bubble formation and disruption of the continuous column.
  \item \textbf{Air ingestion and two-phase breakdown:} air is entrained at the inlet or outlet, creating large bubbles that break the column or change the effective density and viscosity enough to invalidate the model.
\end{enumerate}

A rough estimate of the minimum pressure at the apex in the primed state can be obtained from a Bernoulli-like relation including head losses.
If this estimate falls below the vapor pressure of the liquid, cavitation is expected and the siphon will not sustain a continuous column for long.

\section{Limits of the Model}
\label{sec:limits}

The model developed so far rests on several key assumptions.
Here we summarize the most important ones and indicate when they are expected to fail.

\subsection{Viscous Regime and Reynolds Number}

We have implicitly assumed that the friction factor $f$ in the Darcy--Weisbach formula provides an adequate description of viscous losses.
This requires that the flow be well described by the usual laminar or turbulent pipe-flow correlations, which depend on the Reynolds number
\begin{equation}
  \mathrm{Re} = \frac{\rho v D}{\mu},
\end{equation}
where $\rho$ and $\mu$ are the liquid density and dynamic viscosity, respectively.
At very low Reynolds numbers or in strongly unsteady regimes, standard correlations may be inaccurate, and the effective friction factor may deviate from simple textbook values.

\subsection{Surface Tension and Capillarity}

We have neglected surface-tension effects and treated interfaces as flat on the scale of the tube diameter.
This is valid when gravitational and inertial forces dominate over capillarity.
A useful diagnostic is the Bond number,
\begin{equation}
  \mathrm{Bo} = \frac{\rho g D^2}{\sigma},
\end{equation}
where $\sigma$ is the surface tension.
When $\mathrm{Bo} \gg 1$, surface tension is relatively unimportant; when $\mathrm{Bo} \lesssim 1$, capillary effects can control priming, hysteresis, and the stability of small bubbles.
In very thin tubes, capillary siphons require a different modelling approach.

\subsection{Cavitation and Maximum Crest Height}

The maximum sustainable crest height is constrained by the requirement that the liquid pressure remain above its vapor pressure.
A crude estimate for the pressure at the apex is
\begin{equation}
  p_D \approx p_{\text{atm}} - \rho g (h_D - H) - \rho g h_{\ell,\text{up}},
\end{equation}
where $h_{\ell,\text{up}}$ is the total head loss between the free surface and the apex.
If $p_D$ approaches the vapor pressure $p_{\text{vap}}$, vapor bubbles form and the siphon can break.
This mechanism sets an approximate upper bound on $h_D - H$ even in large-diameter tubes with negligible surface tension.

\subsection{Two-Phase Flow and Air Entrainment}

Our 1D model is strictly single-phase.
In reality, air can be entrained at the inlet, trapped in bends, or ingested at the outlet.
When the volume fraction of air becomes significant, the effective density and viscosity of the mixture change, and the assumptions underlying the Darcy--Weisbach and minor-loss models become questionable.
Severe entrainment can lead to intermittent or slug flow and to failure of the siphon.

\subsection{Compressibility and Water Hammer}

We neglect fluid compressibility and elastic deformation of the tube.
These approximations are valid when flow changes are slow enough that pressure waves are not important.
Rapid start-up or sudden blockage, especially in long or rigid tubes, can generate water-hammer phenomena that require a more detailed transient analysis.

\subsection{Air Resistance}

Finally, we have neglected the aerodynamic drag of the jet and any pressure variations in the ambient air.
For typical laboratory scales, the characteristic dynamic pressure in the liquid, $\rho v^2$, is orders of magnitude larger than that associated with the surrounding air, $\rho_{\text{air}} v^2$, so air resistance can safely be ignored compared to hydrostatic and frictional effects.

\section{Parameter Sensitivity and Intuition}
\label{sec:parameters}

The framework above allows us to interpret how geometric and physical parameters influence the maximum attainable crest height and the likelihood of successful priming.

\begin{itemize}
  \item \textbf{Depth of the inlet and intermediate points (e.g.\ point C):}
  Increasing the submergence depth of the inlet or of an intermediate point below the free surface increases the available hydrostatic head.
  This directly lowers the required initial velocity $v_0$ in the priming criterion \eqref{eq:priming-criterion} and enables higher crests, all else equal.

  \item \textbf{Crest heights:}
  Raising the apex height $h_D$ or the final crest height makes both priming and sustained flow more demanding, because the gravitational potential energy barrier grows and the pressure at the apex decreases.

  \item \textbf{Tube diameter $D$:}
  For a fixed head difference, enlarging $D$ tends to reduce frictional losses (since $h_f \propto L_{\text{tot}} v^2 / D$) but also changes the velocity for a given volumetric flow rate.
  There is a trade-off between adequate velocity for priming and low friction for sustained flow.

  \item \textbf{Tube length and roughness:}
  Longer tubes and rougher walls increase head losses and therefore lower both the steady-state flow rate and the maximum crest height achievable without cavitation.

  \item \textbf{Tank cross-sectional area $A_{\text{tank}}$:}
  A larger tank area leads to a slower decay of $H(t)$ for a given $Q$, prolonging the draining phase and making the flow quasi-steady for longer.
\end{itemize}

These qualitative trends can be made quantitative by nondimensionalizing the governing equations and exploring parameter space numerically.

\section{Suggested Experiments}
\label{sec:experiments}

We briefly outline several experiments that can be performed with simple laboratory equipment to test the predictions and assumptions of the model.

\begin{enumerate}
  \item \textbf{Validation of the loss model:}
  For a fixed tube geometry and large container, measure the steady or quasi-steady flow rate for different outlet height differences.
  Fit an effective friction factor $f$ and minor-loss coefficient $K_m$ and compare with textbook correlations.

  \item \textbf{Priming threshold vs.\ crest height:}
  Using an adjustable apex height, determine the maximum crest height for which momentum-driven priming still occurs.
  Compare with the priming criterion \eqref{eq:priming-criterion}, using measured or estimated values of $v_0$, $f$, and $K_m$.

  \item \textbf{Effect of inlet depth:}
  Repeat the previous experiment while varying the depth of the inlet (or an intermediate point C) below the free surface.
  Quantify how the maximum achievable crest height depends on this depth.

  \item \textbf{Breakdown regimes:}
  Employ very thin tubes to explore regimes where surface tension and capillarity dominate, and increase crest height to approach cavitation limits.
  Document the onset of intermittent flow, bubble formation, and siphon failure.
\end{enumerate}

These experiments provide a pathway to refine the model and to explore the boundary between the continuum, single-phase description and more complex multiphase behavior.

\section{Conclusion}

We have developed a simple one-dimensional model for momentum-driven self-starting siphons, emphasizing the start-up dynamics, the role of geometry, and the principal physical limitations.
The model highlights how initial momentum, available hydrostatic head, and head losses jointly control whether a siphon primes successfully and how high its crest can be.
By making explicit the assumptions and their limits, and by proposing concrete experiments, we aim to provide a framework that is both analytically tractable and physically transparent.

\end{document}
